% Discussion
\section{Discussion}
\label{sec:discussion}

\subsection{Domain Shift and Ranking Inversion}
\label{sec:discussion_shift}

All four detectors experience a severe accuracy drop when transferred from KITTI to Waymo,
with mAP@$[0.50{:}0.95]$ falling by roughly $80$--$90\%$ relative to the in-domain
baseline. This is expected for a direct, no-retraining cross-dataset transfer: KITTI and
Waymo differ in camera calibration, sensor position, scene composition, object-scale
distribution, and annotation density. More consequential is the inversion of the detector
ranking across domains. YOLOv8s, the strongest in-domain detector, shows the lowest
generalization ratio ($0.095$), while RetinaNet, the weakest in-domain detector, generalizes
best ($0.180$). This inversion is a candidate indicator of overfitting to KITTI-specific
statistics for the higher-capacity detectors, and it demonstrates that the in-domain
ranking does not transfer to the out-of-domain setting.

The pattern is consistent with the architecture families. The real-time CNN and transformer
detectors (YOLOv8s and RT-DETR) achieve the highest in-domain accuracy but the strongest
relative degradation, whereas the two Torchvision detectors trained with a conservative,
fixed augmentation policy degrade less in relative terms. This does not establish a causal
mechanism, but it suggests that the choice of training regime and architecture interacts
with domain shift in a way that in-domain validation alone cannot reveal.

\subsection{Small Objects and Vulnerable Road Users}
\label{sec:discussion_safety}

The safety-oriented analysis shows that the dominant failure mode is missed detection of
small objects, and that this is far more severe on Waymo than on KITTI. On KITTI,
small-object recall is $0.88$--$0.96$ across detectors, whereas on Waymo it is
$0.03$--$0.13$. Consequently, the pedestrian-plus-cyclist false-negative rate on Waymo
reaches $0.81$--$0.93$ for all detectors. This result is invisible in average mAP: the
detector with the highest Waymo mAP is not the most safety-reliable, because average
precision is dominated by large, easily detected vehicles. For intelligent-vehicle
perception, small-object and pedestrian/cyclist recall should therefore be treated as
first-class deployment metrics alongside mAP, and safety-critical evaluation should be
reported at a fixed operating point rather than only as an AP curve.

\subsection{Deployment Implications}
\label{sec:discussion_deployment}

The efficiency results reinforce the need for a joint accuracy--efficiency--safety
assessment. YOLOv8s provides an order-of-magnitude latency advantage and the best in-domain
accuracy but the least external stability and safety reliability, making it attractive only
when latency is the binding constraint and in-domain statistics match deployment. RetinaNet
offers the best external generalization and safety-oriented recall but at substantially
higher latency and model size, making it attractive when out-of-domain reliability is
prioritized over speed. RT-DETR-L suffers from a large number of low-confidence
over-detections at the fixed operating point, and Faster R-CNN is the slowest detector
despite mid-range accuracy. The practical implication is that detector selection must be
based on external validation and safety metrics rather than KITTI mAP alone, and that the
choice depends on the deployment priority: latency, in-domain accuracy, or out-of-domain
safety.

\subsection{Interpretation and Limitations of the Comparison}
\label{sec:discussion_interpretation}

These findings are descriptive and correlational: the study quantifies how, where, and why
the four detectors fail under a controlled target-free transfer, but it does not isolate
the causal contribution of architecture, capacity, or training regime. The relative
stability of RetinaNet, for instance, could reflect its conservative training setup as much
as its architecture. Nevertheless, the results provide a reproducible evidence base for
deployment-oriented perception and a concrete motivation for domain-robust training or
adaptation strategies~\cite{chen2018domainadaptive,li2022crossdomainadaptive}, which the
target-free setting deliberately leaves out.
